\chapter{Lenin's Last Days}

The process of economic recovery inaugurated by NEP was overshadowed in 1922 by the onset of Lenin's prolonged and fatal illness. In May 1922 he had a stroke which incapacitated him for many weeks. He returned to work and delivered several speeches during the autumn. But his physical powers were manifestly under strain. On December 12, on medical advice, he withdrew to his apartment in the Kremlin, and there four days later he had a second and severer stroke which permanently paralysed his right side. For the next three months, physical incapacity did not affect his mental faculties; and, though none of the other leaders were apparently allowed to see him, he continued to dictate notes and articles on party affairs. These included the famous ``testament'' of December 25 with its postscript of January 4, 1923. But on March 9, 1923, a third stroke deprived him of speech; and, though he lived for ten months longer, he never worked again. 

After his third stroke hopes of Lenin's eventual recovery gradually faded. The question of the succession loomed ahead, and clouded every other issue. The tightening of party discipline at the tenth party congress in March 1921 had been followed by a party purge, and was further emphasized at the eleventh congress a year later, when 22 dissidents, most of them members of the former Workers' Opposition, were censured, and two of their five leaders expelled from the party; Lenin had asked for the expulsion of all five. This fresh crisis called for a further strengthening of the party machine. The three co-equal secretaries of the party central committee appointed in 1920 (see p. 42 above) had proved ineffective, and were removed from office. On April 4, 1922, a few days after the eleventh congress, it was announced that Stalin had been appointed general secretary, with Molotov and Kuibyshev as secretaries. Nobody found the announcement particularly significant. Stalin was known as a hard-working, efficient and loyal party official. 
% 注：去除了 ".party" 和 "‘‘committee" 的噪点。

When Lenin returned to work after his first stroke, he was evidently alarmed by the way in which Stalin had patiently built up both the power and authority of his office, and his own personal standing; he was now for the first time a leading figure in the party. Lenin did not like either of these developments. He was much preoccupied at this time with the growth of bureaucracy in the state and in the party; and he became acutely mistrustful of Stalin's personality. The testament was dictated, a few days after the second stroke which put in doubt his chances of recovery, in a mood of anxious foreboding. Lenin began with the danger of a split between ``the two classes''---the proletariat and the peasantry---on whose alliance the party rested. This he dismissed as remote. The split which he envisaged as a threat for the ``near future'' was between members of the central committee; and the relation between Stalin and Trotsky was ``a big half of the danger of that split''. Stalin had ``concentrated an enormous power in his hands'', and did not ``always know how to use that power with sufficient caution''. Trotsky, though ``the most able man in the present central committee'', showed ``too far-reaching self-confidence and a disposition to be too much attracted by the purely administrative side of affairs''. Other leading members of the committee did not escape criticism. The hesitation of Zinoviev and Kamenev at the critical moment in October 1917 was recalled; this was ``not, of course, accidental, but . . . ought as little to be used against them personally as the non-Bolshevism of Trotsky''. Bukharin, though ``the biggest and most valuable theoretician in the party'' and ``a favourite of the whole party'', had never fully understood the dialectic, and his views could ``only with the very greatest doubt be regarded as fully Marxist''. It was an unexpected verdict on the man whose ABC of Communism, written in conjunction with Preobrazhensky, and Theory of Historical Materialism were still widely circulated party textbooks. But, however perspicacious Lenin's diagnosis of the shortcomings of his colleagues, the only cure prescribed in the testament was a proposal to increase the membership of the central committee to 50 or 100; and this was unlikely to touch the root of the problem. 
% 注：此处修复了两处极其严重的多栏句子互相穿插的问题，恢复了 "built up both the power..." 和 "The testament was dictated..." 相关的正确语序，并清除了大量如 "ilbuilt", "llhis", "tof a split", "nr" 的乱码。

Lenin's attention had been drawn in the autumn of 1922 to what was happening in Georgia, where procedures for the incorporation of the Georgian republic into the USSR had encountered stiff resistance from the Georgian party committee. A commission headed by Dzerzhinsky visited Georgia in September, and returned to Moscow with the two dissident leaders. At this point Lenin intervened, overruled Stalin who was in charge of the question, and believed himself to have secured a compromise. But he did not follow up the matter, and relations with the Georgians again became embittered. Ordzhonikidze now visited Tiflis and, after a violent struggle, dismissed the rebel leaders, and forced the committee to accept Stalin's proposals. A few days after dictating the testament, Lenin, under what impulse is not clear, returned to the Georgian question. He dictated a memorandum in which he confessed himself ``seriously to blame before the workers of Russia'' for having failed to intervene effectively at an earlier stage. He denounced the recent proceedings as an example of ``Great-Russian chauvinism'', referred to Stalin's ``hastiness and administrative impulsiveness'', and severely censured him, Dzerzhinsky and Ordzhonikidze by name. Then, on January 4, 1923, Lenin's mistrust of Stalin again burst out, and he added a postscript to the testament. Stalin, he now said, was ``too rude'', and should be replaced as general secretary by someone ``more patient, more loyal, more polite, and more attentive to comrades, less capricious etc.''; and as the motive for this recommendation he once more cited the danger of a split, and ``the relation between Stalin and Trotsky''. Finally, early in March, after an occasion on which Stalin was said to have insulted Krupskaya (who had presumably refused him access to Lenin), Lenin wrote a letter to Stalin breaking off ``comradely relations''. Three days later came the third stroke which ended Lenin's active life. 

The approach of the twelfth party congress, which met on April 17, 1923, was a source of embarrassment. Who was to don the mantle of leadership which Lenin had worn without challenge at previous congresses? Lenin's eventual recovery was not yet despaired of. But even an interim choice might pre-judge the future succession. Trotsky, a newcomer to the party with a record of past dissent, owed his commanding position since 1917 to Lenin's unfailing support. Deprived of this, he was an isolated figure, and did not and could not aspire to lead the party. He was regarded with jealous dislike by his immediate colleagues, whom he treated with a certain haughtiness; and his past advocacy of the militarization of labour made him suspect in trade union circles. The three other most prominent leaders---Zinoviev, Kamenev and Stalin---drew together in a determination to block any aggrandizement of Trotsky's role. In this provisional triumvirate Stalin was the junior partner; and he was keenly conscious of the need to live down Lenin's personal hostility, which was probably by this time known to the other leaders, if not to the rank and file of the party. Kamenev had more intelligence than force of character. Zinoviev, weak, vain and ambitious, was only too eager to occupy the empty throne. He presided and spoke at the congress in terms of fulsome subservience to the authority of the absent leader, contriving at the same time to suggest that he was the authorized exponent of Lenin's wisdom. Stalin, by way of contrast, assumed a role of calculated modesty. Claiming nothing for himself, he repeatedly referred to Lenin as his ``teacher'', whose every word he studied and sought to interpret aright. Speaking on organization, he repeated Lenin's strictures against bureaucracy, hypocritically ignoring the fact that these shafts had been aimed in large part at himself. In his report on the national question he emphatically endorsed Lenin's attacks on ``Great-Russian chauvinism'', and gently exculpated himself from the charge of ``hastiness''. Trotsky, clearly anxious to avoid any direct confrontation, absented himself from the debate on the national question. His role at the congress was confined to the presentation of a weighty report on the economic situation, which stated the case for industry and for the ``single economic plan'', but did not directly assail current policies. Latent disagreement with Zinoviev was carefully concealed. 
% 注：修复了多栏穿插断层 ("I' position since 1917... Deprived recovery was not yet despaired of...") 并恢复了正常的英语叙述逻辑。

Throughout the summer of 1923 personal animosities simmered beneath the surface while the economic crisis mounted, and hopes of Lenin's recovery gradually faded away. Though Trotsky was not a candidate for the formal role of leadership, his powerful personality, his record in the civil war, his cogency in argument, and his brilliant oratorical gifts won him widespread popularity in the rank and file of the party, and made him a formidable adversary in any debate on policy. The triumvirate of Zinoviev, Kamenev and Stalin had successfully conspired at the party congress in April to block his advance. They now decided that the moment had come to crush him. The campaign was launched with the utmost caution---partly because Zinoviev and Stalin already, perhaps, did not fully trust one another. 
% 注：删除了 "j Throughout" 中的乱码。

The provocation came from Trotsky's letter of October 8, 1923 (see p. 57 above), which, after its caustic criticism of current economic policies, launched into an attack on the ``incorrect and unhealthy regime in the party''. Nomination had replaced election in appointments to key posts in the party organization; and appointments went to those committed to the maintenance of the existing regime. A ``secretarial apparatus created from above'' had gathered all the threads into its hands, and rendered participation by the rank and file ``illusory''. The letter concluded by demanding that ``secretarial bureaucratism'' should be replaced by ``party democracy''. Coming from a member of the Politburo, it was a formidable indictment, and its cutting edge was pointed unmistakably at Stalin. A few days later the ``platform of the 46'' deplored the rift between the ``secretarial hierarchy'' and ordinary members of the party. The ``dictatorship within the party'' which silenced all criticism was traced back to the emergency decisions of the tenth party congress in March 1921; this regime had ``outlived itself''. The triumvirate could not ignore this open challenge to its authority. 
% 注：去除了 "ithe", "Ition", "I the party", "jmitted", "/apparatus", "finto" 等众多的垂直干扰乱码。

It was at this moment, by a strange fatality, that Trotsky succumbed to the first attack of an intermittent and undiagnosed fever, which continued to afflict him at intervals during the next two or three years. On October 25 the party central committee, in Trotsky's absence through illness, passed a resolution condemning his letter of October 8 as ``a profound political error'' which had ``served as the signal for a fractional grouping (the platform of the 46)''. Throughout November animated discussion of the economic and political issues in the columns of Pravda provoked no intervention either by Trotsky or by the triumvirate. Trotsky's persistent indisposition condemned him to a passive role. But early in December he held parleys with the three leaders, which resulted in an agreed resolution of the Politburo of December 5, 1923. The tactics of the triumvirate were to make maximum concessions to Trotsky on points of principle in order to isolate him from the opposition. The resolution spoke of ``the unique importance of Gosplan'', of the danger of ``bureaucratization'' and ``the degeneration under NEP of a section of party workers'', and of the need for more ``workers' democracy''. The existing preponderance of ``non-proletarian'' elements in the party was to be remedied by ``an influx of new cadres of industrial workers''; this was regarded as a guarantee of ``party democracy''. But the earlier resolution of October 25, in which the party central committee had condemned both Trotsky's letter of October 8 and the platform of the 46, was specifically re-affirmed, so that Trotsky appeared both to renounce his own previous stand, and to acquiesce in the condemnation of those who had come out in his support. Trotsky none the less regarded it as a victory for his principles. 

So artificial a compromise could not last. Three days later Trotsky, still unable to appear in public, expounded his interpretation of the resolution in an open letter which was read at party meetings and published in Pravda. He criticized ``conservatively minded comrades who are inclined to over-rate the role of the machine and under-rate the independence of the party''. He cited the German social-democrats before 1914 as an example of an ``old guard'' which had lapsed into ``opportunism'', and he appealed to the rising generation, which ``reacts most sharply against party bureaucracy''. In a postscript he referred to ``the dangers of NEP'', closely connected with ``the protracted character of the international revolution''. The triumvirate still hesitated. At a meeting of the Moscow party organization on December 11, several supporters of Trotsky spoke, including Preobrazhensky and Radek; and Zinoviev and Kamenev, while condemning the opposition, handled Trotsky with cautious politeness. 
% 注：修复了断字拼写错误 ("theprotracted" 改为 "the protracted")。

A few days later all inhibitions were removed, and the triumvirate had decided to treat Trotsky's open letter as a declaration of war. On December 15, Stalin in an article in Pravda launched a full-scale attack on the opposition, ending with bitter personal gibes against Trotsky. This was the signal for a campaign of vituperation in a series of speeches and articles by Zinoviev, who seems to have coined the term ``Trotskyism'', Kamenev, Bukharin and lesser party figures. Articles favourable to the opposition no longer appeared in Pravda. Students demonstrated for the opposition; and a purge of the central committee of the Komsomol was conducted in order to bring that organization to heel. But at party meetings in Moscow and Petrograd only a small minority of workers spoke or voted against the official line. Trotsky's earlier advocacy of the militarization of labour made it difficult for him to appear as a champion of the workers' cause. The growing power of the party organization, the lack of any positive or popularly presented alternative programme, fear of victimization in a period of widespread unemployment, the weakness in numbers and in radical traditions of the Russian working class, all contributed to the rout of the opposition. A protest by Trotsky, Radek and Pyatakov against the discriminatory attitude of Pravda produced a reply from the party control commission that ``the organ of the central committee is obliged to carry out the perfectly definite line of the central committee''. The decision was final and absolute. Thereafter Pravda spoke exclusively with the official voice of the central organs of the party. 
% 注：去除了 "jany", "[fear", "Ithe", "'committee", "/of", "'Thereafter" 等诸多乱码。

The process of the personal denigration of Trotsky quickly gathered strength. At a session of IKKI early in January 1924 Zinoviev delivered a further uninhibited attack on his character, his party record and his opinions. Trotsky, dogged by illness, gave up the unequal struggle, and departed on medical advice for the Caucasus in the middle of January 1924. A few days later a party conference by an overwhelming majority (the delegates had, no doubt, been carefully sifted) condemned the opposition, and held Trotsky personally responsible for the campaign against the party leaders. These events immediately preceded the death of Lenin, which took place on January 21, 1924.